\section{Resultados}
\label{sec:resultados}

Aplicando o algoritmo da Seção~\ref{sec:algoritmo} aos oito casos de teste, com a regra do
pulo da Seção~\ref{sec:regra}, obtêm-se os resultados da Tabela~\ref{tab:resultados}.

\begin{table}[ht]
  \centering
  \caption{Resultados dos oito casos de teste. A coluna ``Casas retiradas'' conta as casas que
  saíram da fila até a retirada de S, e o tempo é o melhor de 15 execuções.}
  \label{tab:resultados}
  \small
  \begin{tabular}{|c|c|c|r|c|r|}
    \hline
    \textbf{Caso} & \textbf{$N$} & \textbf{Movimentos} & \textbf{Casas retiradas}
      & \textbf{\% de $N^2$} & \textbf{Tempo (ms)}\\
    \hline
    caso40 & 40 & 4 & 425 & 26,6 & 0,02\\
    \hline
    caso80 & 80 & 6 & 4.483 & 70,0 & 0,40\\
    \hline
    caso100 & 100 & 8 & 8.596 & 86,0 & 0,78\\
    \hline
    caso150 & 150 & 11 & 20.863 & 92,7 & 1,92\\
    \hline
    caso200 & 200 & 9 & 15.483 & 38,7 & 1,50\\
    \hline
    caso400 & 400 & 19 & 92.523 & 57,8 & 8,79\\
    \hline
    caso800 & 800 & 43 & 597.282 & 93,3 & 57,24\\
    \hline
    caso1500 & 1500 & 34 & 375.230 & 16,7 & 36,30\\
    \hline
  \end{tabular}
\end{table}

Nenhum dos oito casos é impossível: em todos, S é alcançável a partir de C. O número de
movimentos vai de 4 a 43, e o tabuleiro de maior lado, $1500 \times 1500$, foi resolvido em
36~ms. O caso mais demorado é o 800, com 57~ms, e a razão disso, que não é o tamanho, está na
Seção~\ref{sec:casos-reais}.

\subsection{Como os resultados foram conferidos}
\label{sec:conferencia}

Para não confiar em um único programa, os resultados foram conferidos de várias formas:

\begin{itemize}
  \item \textbf{Exemplo do enunciado.} O algoritmo dá 3 pulos, como o enunciado informa, e uma
    busca exaustiva confirma que 2 pulos não bastam.
  \item \textbf{Segunda implementação.} A busca foi escrita de novo, separadamente (casas
    guardadas como números inteiros e fila em vetor), para que um erro do primeiro programa não
    se repetisse igual na conferência. As duas dão a mesma resposta nos oito casos.
  \item \textbf{Algoritmo de Dijkstra.} Um método diferente, com fila de prioridade, dá a
    mesma distância nos casos até $400 \times 400$.
  \item \textbf{Caminho legal.} Reconstruiu-se o caminho mínimo de cada caso e conferiu-se, pulo
    a pulo, que todos obedecem à regra do pulo.
  \item \textbf{Tabuleiros pequenos aleatórios.} Em 389 tabuleiros sorteados (com semente
    fixa), a distância encontrada é atingível e uma busca exaustiva em profundidade limitada
    não acha caminho menor.
  \item \textbf{Casos de borda.} Tabuleiros construídos à mão testam o pulo padrão, um pulo
    que só fecha dando a volta no toro e o dígito 9 em um tabuleiro $7 \times 7$ (pernas
    maiores que o lado, o que testa o módulo). A saída ``impossível'' foi exercitada com um
    tabuleiro sem solução, e a busca independente também não achou caminho.
\end{itemize}

Essas conferências mostram que o programa implementa corretamente a regra do pulo da
Seção~\ref{sec:regra}. Como as implementações independentes seguem a mesma regra, elas testam a
implementação, e não a interpretação do enunciado.
